\scanpage{351}
\mspage{\deleted{\unsure{89|80}}90}
% archive sheet number: 350
\deleted{\unsure{deux|aux} valeurs de $F$ en dehors de $I$ (donc des homotopies qui soient l'identité sur les sommets $s\in I_\tau\setminus I$, et de plus de \inserted{$p={}$}$0_I$ \deleted{\unsure{$*$}} et \inserted{$q={}$}$1_I$), en un $F'$ qui, sur $I$, n'aura qu'une seule flèche \inserted{de tr.} critique, et celle-ci $=$ \unsure{plus|p. ex.} au début, égale à $\underline{u}_{i_\alpha}$ (si on y tient).}

\deleted{Il faut cependant traiter à part le cas $p=0,q=n$. Alors les $\underline{u}_{i_\alpha}$ sont \uline{toutes} de $=$ signe, Et il s'agit d'homotoper $F$ à un foncteur $F'$ qui n'a qu'une seule flèche critique. \deleted{Si $\tau$ est \unsure{bien un} $\Delta_\varepsilon^n$, Si les $\underline{u}_{i_\alpha}$ ne \unsure{sont pas tous}, alors par l'hyp. de} récurrence on est ramené au cas, d'abord où sur chaque intervalle critique de $I_\tau$, il y a au plus une flèche critique. (il suffit de l'assurer sur le dernier intervalle critique, contenant la flèche $F$-critique $\underline{u}_n$), puis au cas où il n'y a que deux flèches critiques, \deleted{\unsure{et}} savoir $u_1$ et $u_n$. Ensuite, \deleted{on se ramène \unsure{encore} par l'hyp. de récurrence, en \unsure{ramenant} \unsure{un à à}} Utilisant l'hyp. de réc. pour $I=[1,n-1]\subset I_\tau$, etc \unsure{si on})}

1) On veut assurer a). Pour ceci, on utilisera l'hyp. de récurrence, pour \deleted{\unsure{grouper ensemble toutes les}} \inserted{remplacer un groupe \unsure{de au moins|de moins} $1$} flèches critiques consécutives de $=$ signe, par une seule flèche critique. \unsure{Ladite|cette} méthode ne foire que si les flèches critiques sont toutes de $=$ signe \inserted{$\varepsilon$} et \unsure{incluent|enclavent} $\underline{u}_1$ et $\underline{u}_n$. Mais l'hyp. de récurrence appliquée à $I=[1,n-1]\subset I_\tau$ nous ramène au cas où il n'y a que deux flèches critiques $u_1$ et $u_n$. Si le premier intervalle critique n'est pas réduit à $\Delta_\varepsilon$ i.e. est de longueur $2$, alors l'hyp. de récurrence appliquée à $\unsure{J|I}=[1,\unsure{n}]\subset I_\tau$ nous ramène au cas où les deux flèches critiques sont $\underline{u}_1$ et $\underline{u}_2$, \deleted{\unsure{donc alors}} \inserted{et on est}
